\subsection{The actual polarized model at a coefficient curve}
\label{sec:selmer-polarized-model}

Retain the original ordinary pair \((D,\psi)\) over \(B\), its
normalized point \(\rho:G_0\to\mathrm{GL}_n(k[[t]])\), the
surjection \(h:G\to G_0\), and all prescribed ordinary and auxiliary
equations.  Write \(Z=\mathrm{pairGens}\cup\mathrm{qGens}\) and
\(I=D.\mathrm{chBar}(Z)=\mathrm{OoPt.ker}\) for the exact original
ideal.  It is the intersection of
\((D.\mathrm{CH},Z,\mathfrak m_B^j,[u]-1\ (u\in U))\) over all
\(j\) and open normal \(U\), already including coefficient and
group closure.  Put
\[
 F=k((t)),\qquad R=\widehat{B_{\ker f}},
 \qquad f:B\longrightarrow k[[t]].
\]
The specified curve variable maps to \(t\), and \(\mathcal O\to k\)
is surjective, so \(f\) is surjective and \(R\) has residue field
\(F\).  The actual normalized point has polarization matrix \(I_n\)
and group span \(M_n(F)\), or \(M_a(F)\times M_{n-a}(F)\).
Assume \(\ell=\operatorname{char}k\ne2\), as in the source.

\begin{proposition}[The full polarized equation model]
\label{prop:selmer-polarized-model}
There is a finite commutative \(R\)-algebra
\(\mathcal A=R\oplus N\), with \(N^2\subseteq\mathfrak m_R\), and
an actual representation \(v:G\to\mathrm{GL}_n(\mathcal A)\).
It has the original determinant and polarization, kills the
original ideal \(I\), and its \(R\)-linear span consists exactly
of the matrices with even-block entries in \(R\) and cross-block
entries in \(N\).  It retains the original residual matrix entries
and finite-\(B\)-lattice continuity.  If the original equation
quotient at a prime \(Q\subset R\) is nonsemisimple, then
\[
                         \operatorname{Ann}_R N\subseteq Q.
\]
When the normalized point is irreducible, \(N=0\).
Thus all fields of the source's \(\mathrm{URep}/\mathrm{S6}\)
construction are supplied.
\end{proposition}

\begin{proof}
\emph{The exact equation algebra.}
The supplied limiting datum makes \(G.\Gamma\) topologically
finitely generated.  Profinite Schreier gives the same property for
the actual equation group \(G.\mathrm G_0=\ker\eta\), an open
index-two subgroup, and for its open subgroups.  Thus continuous
cocycles into each finite module form a finite set.
Wang--Erickson's Proposition 3.6 gives a
finite continuous Cayley--Hamilton algebra for its determinant,
with the group, quotient and coefficient-adic topologies agreeing
\cite{SelmerModelWE}.  Base change to \(B\) retains these properties.
The original residual characters are distinct.  The
generalized matrix algebra theorem and its full-determinant
version apply in every residue characteristic
\cite[Theorem 2.22]{SelmerModelChenevier}
\cite[Theorem 2.27]{SelmerModelWE}.

The prescribed relation generators \(z\) satisfy
\(\operatorname{tr}D(gz)=0\) for all \(g\), and their closed
two-sided ideal lies in the trace radical.  This quotient retains
the full determinant without division by \(n!\).  Indeed each
primitive corner is \(Be_i\) with \(\operatorname{tr}(e_i)=1\);
a trace-radical ideal \(J\) therefore has \(e_iJe_i=0\).
In the canonical determinant formula, a term containing an
element of \(J\) has a permutation cycle lying in some
\(e_iJe_i\), and vanishes
\cite[Proposition 2.23]{SelmerModelWE}.  This argument survives
every coefficient extension.

The completed group algebra quotient has exactly this algebraic
kernel.  Indeed, in
\(\widehat{B[G]}=\varprojlim_{j,U}(B/\mathfrak m_B^j)[G/U]\)
the closed ideal generated by \(D.\mathrm{CH}\cup Z\) intersects
\(B[G]\) in \(I\), by the displayed intersection formula.
Continuity of the universal polynomial identities identifies its
quotient with the finite Cayley--Hamilton equation algebra above.
The image of \(B[G]\) is dense, but also closed: it is a submodule
of a finite module over the complete Noetherian \(B\).  It is
therefore surjective.  Consequently
\[
 E_0=B[G]/I,\qquad
 E=R\otimes_BE_0=R[G]/(R\cdot I),
\]
where \(R\cdot I\) is the algebraic two-sided span of the image of
the already-closed original \(I\), with no further closure or
saturation.  Right exactness gives the second equality without
flatness; this is precisely the source's \(\mathrm{Hat.ker}\).
The full determinant and involution are retained:
\[
                     [g]^*=\varepsilon(g)^{1-n}[cgc]^{-1}.
\]
The finite \(B\)-module \(E_0\) is continuous.  The finite complete
\(R\)-algebra \(E\) has the prescribed semisimple residual quotient.
The full fibre \(E/\mathfrak m_RE\) may retain a nilpotent radical
and \(N/\mathfrak m_RN\); neither is discarded.

\emph{The polarization fixes the product sign.}
Since the normalized polarization matrix is \(I_n\), the
residual involution is ordinary transpose on each block.
It neither exchanges the blocks nor introduces an alternating
sign.  Lift its primitive diagonal idempotents to self-adjoint
orthogonal \(e_i\in E\).  Explicitly, average a lift with its
star using \(1/2\), then iterate
\[
                     u\longmapsto u-(u^2-u)(2u-1)^{-1}
\]
in the radical-adic topology, proceeding in complementary
corners.  The defect squares, the iterates remain self-adjoint,
and the radical-adic and coefficient-adic topologies agree.
The rank-one Cayley--Hamilton corner is \(e_iEe_i=Re_i\),
with trace one.

In each block fix a reference \(r\).  A lift \(u_i\in e_iEe_r\)
of \(E_{ir}\) satisfies
\[
 u_i^*u_i=c_ie_r,\qquad u_iu_i^*=c_ie_i,\qquad c_i\in1+\mathfrak m_R.
\]
Cyclic trace identifies the two scalars.  Hensel lifting the
square root of \(c_i\) normalizes \(g_i=u_ic_i^{-1/2}\).
The units \(E_{ij}=g_ig_j^*\) satisfy \(E_{ij}^*=E_{ji}\).
No coefficient-field extension or division by a block rank is used.

If there is one block, this already gives \(E=M_n(R)\) and
\(N=0\).  For two blocks, let \(e,f\) be their reference
primitive idempotents and set \(N=eEf\).
Identify \(fEe\) with \(N\) by star.  Define
\[
                         mn^*=b(m,n)e.
\]
Applying star proves \(b(m,n)=b(n,m)\).  Cyclic trace, using
\(\operatorname{tr}(e)=\operatorname{tr}(f)=1\), gives
\(n^*m=b(m,n)f\).  Associativity in \(E\) then gives the crucial
third identity
\[
                         b(m,n)p=b(n,p)m\qquad(m,n,p\in N).
\]
Thus
\[
 (r,m)(s,n)=\bigl(rs+b(m,n),\,rn+sm\bigr)
\]
makes \(R\oplus N\) a commutative associative algebra.
The last identity checks associativity for three odd factors.
The residual block algebra kills cross corners, so
\(b(N,N)\subseteq\mathfrak m_R\).
This keeps all annihilator torsion in \(N\).

The matrix units identify
\[
 E=
 \begin{pmatrix}
 M_a(R)&M_{a,n-a}(N)\\
 M_{n-a,a}(N)&M_{n-a}(R)
 \end{pmatrix}
 \ \subseteq M_n(R\oplus N).
\]
The involution is transpose.  The determinant cycle formula is
unchanged under this embedding, so the entire determinant law is
the original one, including in rank three at \(\ell=3\).
The group map gives \(v\), and its linear span is exactly this
displayed algebra.  The polarization is precisely
\[
                    v(g)^Tv(cgc)=\varepsilon(g)^{1-n}I_n.
\]

\emph{The prime and continuity fields.}
If \(c\in\operatorname{Ann}_R N\setminus Q\), every field-valued
map of kernel \(Q\) kills \(N\).  The original equation quotient,
using the exact kernel equality above, becomes
\(M_a(\kappa)\times M_{n-a}(\kappa)\), a semisimple algebra.
This contradicts the original nonsemisimplicity witness and
proves the annihilator inclusion for all such witnesses.

Let \(P_0\) be the image of \(E_0\) in \(M_n(\mathcal A)\).
It is finite over the original \(B\), contains every \(v(g)\),
and is closed under multiplication.  Continuity gives, for
each \(j\), an open normal \(V\subset G\) with
\(v(u)-1\in\mathfrak m_B^jP_0\) for \(u\in V\).
Multiplication by \(v(g)\) therefore gives exactly
\(v(gu)-v(g)\in\mathfrak m_B^jP_0\).
The original equations and every remaining model field now follow
from the displayed construction \cite{SelmerModelSource}.
\end{proof}

Let \(\mathfrak r=(b(N,N))\).  If
\(\mathfrak r\subseteq Q\), the trace pairing modulo \(Q\)
factors through \(M_a\times M_{n-a}\), of dimension
\(a^2+(n-a)^2<n^2\).  Hence the stable discriminant condition
forces \(\mathfrak r\nsubseteq Q\), exactly as required by the
local incidence theorem.  This supplies the actual even--odd
algebra in that theorem; the remaining tangent/presentation,
patched support and arithmetic seed obligations are separate.
